\chapter{Infrastructure: Kaggle and RunPod}\label{app:infra}

\section{Where the computing happens}\label{sec:infra-where}
Almost none of the computation in this project runs on the computer that holds the repository. From
15 September 2026 onwards we followed an explicit rule: nothing is downloaded to the local machine,
and no local bandwidth is spent. Code is written locally, shipped to remote machines and run there,
and only small artefacts, such as reports, CSV files, figures and PDF documents, come back. We used two
kinds of remote machine: Kaggle kernels, which are free within a weekly quota, and GPUs rented from
RunPod, which are paid for by the hour.

\section{Kaggle kernels}\label{sec:infra-kaggle}
A Kaggle \keyterm{kernel} (which Kaggle now also calls a notebook) is a script that Kaggle runs on its
own hardware, with the competition data and any chosen datasets mounted read-only under
\file{/kaggle/input}. The following facts about kernels shaped how we worked.
\begin{itemize}[leftmargin=*,itemsep=2pt]
\item \textbf{Hardware.} Kaggle offers kernels with CPUs only, with one or two NVIDIA T4 GPUs (16 GB
each), or with a TPU v5e with eight cores. During the project our account had about 30 GPU hours and
20 TPU hours per week, one TPU session at a time, and a limit of 12 hours per session.
\item \textbf{Outputs.} Only the files written under \file{/kaggle/working} survive a session. An early
experiment lost 1.5 hours of training because it wrote its checkpoints to \file{/kaggle/temp}.
\item \textbf{Metadata.} Each kernel folder contains a file \file{kernel-metadata.json}, which names
the script, the datasets and other kernels to mount, and whether a GPU and internet access are enabled.
Kaggle derives the address of a kernel (its ``slug'') from its \emph{title}, not from the \texttt{id}
field of this file. When the two disagree, a push can land at an address that no later status query
looks at, so we keep them identical.
\item \textbf{Knowing what runs.} The command \texttt{kaggle kernels list --user <user> --sort-by
dateRun} lists the kernels that actually ran. A status query for a single guessed name is not a
reliable check (Section~\ref{sec:lessons-config}).
\item \textbf{Chaining sessions.} A kernel can mount its own earlier output, which in principle allows
twelve-hour sessions to be chained. However, a kernel that is stopped by the time limit is
\emph{cancelled}, and the output of a cancelled kernel cannot be mounted. Any plan that relies on
chaining must be tested on that path.
\item \textbf{CPU and GPU agree.} Inference on Kaggle's CPUs produced a submission file identical, bit
for bit, to the one produced by inference on a GPU (the two files had the same MD5 hash). We therefore
used CPU kernels freely for inference and scoring, and kept the GPU hours for training.
\end{itemize}

\subsection{TPUs}\label{sec:infra-tpu}
Kaggle's TPU trained the dense U-Nets well: exp004 ran 300 epochs at 1024 pixels in 3.8 hours. It used
a single Python process with SPMD sharding across the eight cores, because the usual launcher, which
starts one process per core, failed while the TPU was starting up on Kaggle's image. The XLA compiler
recompiles the model for every new tensor shape, so the batches had to have fixed shapes, and every
call to \texttt{.item()} stalls the pipeline, so we read the losses only once per epoch. YOLO training
never ran on the TPU, because compilation timed out, and the first Mask2Former trained there far too
slowly to finish (Section~\ref{sec:train-first}).

\section{Datasets as a means of transport}\label{sec:infra-datasets}
Private Kaggle datasets carry code to the remote machines and bring weights back. Two properties of
Kaggle's datasets determined how we use them.
\begin{itemize}[leftmargin=*,itemsep=2pt]
\item \textbf{Kaggle unpacks archives.} After an upload, a \file{.tar.gz} file is listed as the files
it contains, so a single archive whose SHA-256 hash a remote machine intended to check no longer
exists. We therefore ship code as a tar archive renamed to \file{.payload}, and the receiving machine
verifies its SHA-256 hash before running anything.
\item \textbf{The listing must be checked.} After every upload we list the dataset's files and their
sizes, and compare them with what we intended to upload, before anything depends on the dataset.
\end{itemize}
On Windows, the version of the Kaggle command-line tool that we used failed on \texttt{-p} arguments
containing forward slashes; running the tool from inside the folder with \texttt{-p .} avoided the
problem. Downloading the logs of a kernel required setting \texttt{PYTHONUTF8=1}, because otherwise a
single non-ASCII character in the log produced an empty file.

\section{RunPod}\label{sec:infra-runpod}
We trained two models on rented GPUs:
\begin{center}\small
\begin{tabular}{@{}llccc@{}}
\toprule
experiment & GPU & per hour & length & cost \\
\midrule
exp097 (YOLOv8l-seg) & L40, 48 GB & \$0.69 & 124 epochs & \$2.74 \\
exp139 (Mask2Former) & H100 NVL, 94 GB & \$2.59 & 50 epochs & about \$4.80 \\
\bottomrule
\end{tabular}
\end{center}
Two early machines for exp097 that failed before training began added \$0.22 to the cost.

A rented machine starts from a standard PyTorch image and runs a bootstrap script. The script downloads
the code payload from a private Kaggle dataset, checks its SHA-256 hash, reproduces a known validation
score to prove that the scorer works on that machine, runs a short preflight that measures speed and
memory, and only then starts training. Checkpoints are uploaded to a private Kaggle dataset, whose
listing the machine verifies before it stops itself. If any command fails, an error handler uploads the
logs and stops the machine, so that a failure costs minutes rather than hours. For the same reason, the
diagnostics that a failed machine uploads, and not the provider's own log, are the place to look for the
cause of a failure (Section~\ref{sec:train-run}).

The Kaggle access key reaches a rented machine only through its environment, when the machine is
created, and it is never written to a file in the repository or to a log. Spending is bounded in two
ways: the trainer projects its cost from its measured speed and stops at a cap, and the hourly price of
every new machine is stated and approved before the machine is created.

\section{Submitting and watching}\label{sec:infra-submit}
A submission is a single command:
\begin{lstlisting}[language=bash,caption={Submitting a CSV file to the competition.},label=lst:infra-submit]
kaggle competitions submit -c filament-segmentation-2026 -f submission.csv -m "exp139 epoch 40 ..."
\end{lstlisting}
The message attached to a submission has a practical use. The script
\file{scripts/watch\_and\_submit.sh} polls kernels, downloads each finished submission file and submits
it, and it decides whether a file has already been submitted by comparing messages. An earlier version
compared file names instead, and since every kernel calls its output \file{submission.csv}, it skipped a
new file because an old submission had the same name. Finally, \file{scripts/alarm.ps1} plays a sound and
shows a window when a long job ends:
\begin{lstlisting}[language=bash,caption={The end-of-job alarm on Windows.},label=lst:infra-alarm]
powershell -ExecutionPolicy Bypass -File scripts\alarm.ps1 -Message "book built" -Seconds 15
\end{lstlisting}
